\begin{solapartat}
\punts{0,1} $T_{1/2} = 7,00\times 10^{8}\ \textit{anys}\cdot\dfrac{(365\times 24\times 3600)\ \mathrm{s}}{1\ \textit{any}} = 2,21\times 10^{16}\ \mathrm{s}$

La constant de desintegració es calcula a partir del període de semidesintegració, que és el temps que ha de passar perquè es redueixi a la meitat la quantitat d'una substància radioactiva.

\punts{0,3}
\[ \left.\begin{aligned} &N_{1/2} = N_0 e^{-\lambda T_{1/2}}\\ &N_{1/2} = \frac{1}{2}N_0\end{aligned}\right\}\lambda = \frac{\ln 2}{T_{1/2}} = 3,14\times 10^{-17}\ \mathrm{s^{-1}} \]

El nombre de nuclis inicials ($N_0$):

\punts{0,2} $N_0 = 1,00\ \mathrm{g}\cdot\dfrac{1\ \text{mol d'U}}{235\ \mathrm{g}}\cdot\dfrac{6,02\times 10^{23}\ \textit{nuclis}}{1\ \mathrm{mol}} = 2,56\times 10^{21}\ \textit{nuclis}$

\punts{0,2} L'activitat d'una substància radioactiva es defineix com:
\[ A(t) = -\frac{dN(t)}{dt} = (-\lambda)\underbrace{(N_0)e^{-\lambda t}}_{N(t)} \]

\punts{0,2} $A_0 = \lambda N_0 = 8,04\times 10^{4}\ \mathrm{Bq}$
\end{solapartat}

\begin{solapartat}
\punts{0,2} $m(t) = m_0 e^{-\lambda t}$

\punts{0,2} $t = 10^{8}\ \textit{anys}\cdot\dfrac{(365\times 24\times 3600)\ \mathrm{s}}{1\ \textit{any}} = 3,15\times 10^{15}\ \mathrm{s}$

\punts{0,6} $m(t = 10^{8}\ \textit{anys}) = 1,000\times e^{-3,14\times 10^{-17}\times 3,15\times 10^{15}} = 0,906\ \mathrm{g}$
\end{solapartat}
